% Mechanical relocation generated from the preservation ledger.
% Existing prose is included byte-for-byte from preserved blocks.

\section{Adequacy and Identification}
\subsection{Full-Data Diagnostics}

Full-data adequacy asks whether fitted means reproduce the observations used in
estimation. Residuals and count-model deviance should be reported overall and
by measurement type, line, direction, stop, period, and vehicle journey where
those labels exist. Systematic grouped residuals indicate misspecification,
data problems, or an inadequate observation model; a small aggregate error can
hide compensating local errors.

Identification is a different question. It concerns sensitivity of the
likelihood to parameter changes and is assessed through response derivatives,
curvature eigenvalues, condition measures, profiles, prior sensitivity, and
replicated fits. A model can fit calibration counts closely while leaving some
parameters weakly identified. A strongly identified but misspecified model can
also fit poorly. The two diagnoses must be reported separately.

For parameter vector $\bm\eta$, local curvature of the negative log posterior
or negative log likelihood summarizes sensitivity near the fitted point. Small
eigenvalues identify locally weak directions; a large condition number warns
that reported parameter precision may be unstable. These local quantities
should be complemented by prior sensitivity, profiles, and repeated fits
because bounded, asymmetric, or multimodal objectives are not characterized by
one Hessian.

The reduced-OD fit result operationalizes these checks without reconstructing
the complete OD vector. It reports the symmetric Hessian eigensystem, numerical
rank, condition estimate, and weak eigenvector loadings using production-basis
labels where available. Negative curvature and numerical singularity produce
warnings rather than being hidden by a covariance calculation. For MAP fits,
the report separates likelihood and prior contributions, gives the penalty of
each raw parameter, standardizes displacement from the prior mean, and warns
when the prior dominates the fitted objective.

Production diagnostics are calculated from the origin--period production
vector, not by materializing all OD cells. They include fitted-to-baseline
ratios, total production, dispersion proximity to its numerical floor, and
extreme production multipliers. These summaries test numerical and scientific
plausibility but do not convert operational boarding or alighting events into
journey productions. Active bounds, extreme multipliers, poor conditioning,
or strong prior dominance must therefore be reported alongside residual
adequacy.

When the objective is to plan additional data collection, these diagnostics
must be evaluated for explicit candidate measurements. Expected information
gain, improvement in conditioning, or reduction in uncertainty can support a
counter-placement decision. Existing grouped residuals alone cannot: they
identify where the current model fits poorly, not necessarily where an
additional count would resolve a weak parameter direction.
